\section{Lubang Hitam Schwarzschild}
\label{sec:lubang-hitam}

\noindent Lubang hitam merupakan salah satu prediksi dari teori relativitas 
umum, yaitu suatu daerah dalam ruang-waktu yang memiliki medan gravitasi sedemikian kuat 
sehingga tidak ada benda apa pun, termasuk cahaya, yang dapat lepas dari dalamnya begitu 
telah melewati suatu batas tertentu \parencite{Frolov1998}. Secara lebih formal, 
lubang hitam dapat didefinisikan sebagai daerah ruang-waktu yang secara kausal terputus 
dari pengamat pada jarak tak hingga, yakni tidak ada sinyal fisis maupun informasi apa 
pun dari daerah tersebut yang dapat mencapai daerah yang jauh \parencite{Chandrasekhar1983}. 
Berbeda dengan objek astrofisika pada umumnya, lubang hitam tidak dicirikan oleh permukaan fisis yang padat, melainkan oleh 
struktur kausal ruang-waktunya sendiri, yang batasnya dikenal sebagai horizon peristiwa.

Horizon peristiwa merupakan permukaan yang membatasi daerah lubang hitam dari daerah 
ruang-waktu di luarnya, di mana segala sesuatu yang melintasinya menuju ke dalam, baik 
materi maupun radiasi, tidak dapat lagi kembali keluar, sehingga daerah yang masih dapat 
berkomunikasi dengan pengamat di kejauhan terpisah secara permanen dari daerah di baliknya 
\parencite{Frolov1998, Wald1984}. Meskipun demikian, horizon peristiwa berbeda dari permukaan fisis biasa seperti dinding 
atau batas suatu objek padat. Seorang pengamat yang melintasi horizon peristiwa tidak akan 
merasakan keganjilan fisis apa pun secara lokal pada saat melintas, karena horizon 
peristiwa bukan merupakan tempat dengan kelengkungan tak hingga, melainkan semata-mata 
penanda batas kausal. Dalam pembahasan pada penelitian ini, horizon 
peristiwa berperan penting sebagai salah satu titik batas domain fisis yang menjadi acuan 
kondisi batas.

Di samping horizon peristiwa, lubang hitam juga dicirikan oleh keberadaan singularitas, 
yaitu di titik pusat lubang hitam tempat kelengkungan ruang-waktu menjadi tak hingga dan 
deskripsi geometri klasik tidak lagi berlaku \parencite{Frolov1998}. Berbeda dengan horizon 
peristiwa yang hanya menandai batas kausal, singularitas merupakan tempat kelengkungan 
yang sesungguhnya tak hingga dan tidak dapat dihilangkan oleh transformasi koordinat 
apa pun.

Berbeda dengan objek astrofisika lain yang memerlukan banyak besaran untuk 
mendeskripsikannya secara lengkap, lubang hitam justru dicirikan oleh jumlah besaran yang 
sangat sedikit. Berdasarkan teorema ketiadaan rambut (\textit{no-hair theorem}), lubang 
hitam sepenuhnya dicirikan hanya oleh tiga besaran, yaitu massa, momentum sudut, dan muatan 
listrik \parencite{Frolov1998, Ruffini1971}. Pada geometri ruang-waktu yang statis dan 
bersimetri bola sebagaimana dideskripsikan oleh solusi Schwarzschild, momentum sudut dan 
muatan listrik bernilai nol, sehingga diperoleh jenis lubang hitam paling sederhana yang 
sepenuhnya dicirikan hanya oleh massanya.

Cara paling langsung untuk memahami struktur geometri lubang hitam Schwarzschild adalah 
dengan meninjau perilaku kerucut cahaya (\textit{light cone}) di sekitarnya 
\parencite{Carroll2004}. Lintasan cahaya radial diperoleh dengan menetapkan $\theta$ dan 
$\varphi$ konstan serta $ds^2 = 0$ pada elemen garis Schwarzschild. Karena $\theta$ dan 
$\varphi$ tidak berubah sepanjang lintasan, perubahan infinitesimalnya $d\theta$ dan 
$d\varphi$ bernilai nol, sehingga elemen garis tereduksi menjadi
\begin{equation}
    0 = -f(r)\,dt^2 + f(r)^{-1}\,dr^2,
\end{equation}
yang menghasilkan laju koordinat cahaya radial
\begin{equation}
    \frac{dt}{dr} = \pm\frac{1}{f(r)}.
    \label{eq:laju-cahaya-radial}
\end{equation}
Pada jarak yang jauh dari lubang hitam ($r \gg 2M$), $f(r) \to 1$ sehingga $dt/dr \to \pm 1$ 
dan kerucut cahaya tampak terbuka lebar seperti pada ruang-waktu datar. Akan tetapi, ketika 
$r$ mendekati $2M$, fungsi metrik $f(r) \to 0$ sehingga $dt/dr \to \pm\infty$, yang berarti 
kerucut cahaya semakin menyempit dan akhirnya menutup sepenuhnya tepat di $r = 2M$ 
\parencite{Carroll2004}. Penutupan kerucut cahaya ini secara fisis menunjukkan bahwa cahaya 
yang berada di permukaan $r = 2M$ tidak lagi dapat merambat menjauh dari lubang hitam, 
melainkan hanya dapat menyusuri permukaan itu sendiri atau jatuh ke dalam.

Secara formal, perilaku ini berkaitan langsung dengan sifat kausal dari permukaan 
$r = \text{konstan}$ itu sendiri. Permukaan horizon peristiwa secara definisi merupakan 
permukaan nol (\textit{null hypersurface}), yaitu permukaan yang hanya dapat dilintasi ke 
satu arah oleh kurva kausal \parencite{Carroll2004}. Untuk menentukan letak permukaan ini, 
setiap permukaan $r = \text{konstan}$ memiliki vektor normal yang diberikan oleh gradien 
fungsi $r$,
\begin{equation}
    \eta_\mu = \partial_\mu r,
\end{equation}
yang secara alamiah berbentuk \textit{one-form} (kovektor) dan selalu tegak lurus terhadap 
permukaan $r = \text{konstan}$ \parencite{Carroll2004}. Sifat kausal permukaan ditentukan 
oleh norma vektor normalnya. Karena $\eta_\mu$ berindeks bawah, normanya dihitung 
menggunakan invers metrik $g^{\mu\nu}$. Mengingat bahwa $\partial_\mu r$ hanya memiliki 
komponen tak nol pada $\mu = r$, perhitungan norma secara langsung menghasilkan
\begin{equation}
    g^{\mu\nu}\eta_\mu\eta_\nu = g^{rr}.
\end{equation}
Ketika $g^{rr} = 0$, vektor normal bersifat \textit{null} sehingga permukaan menjadi 
\textit{null hypersurface} yang hanya dapat ditembus satu arah, yang persis merupakan 
definisi horizon peristiwa \parencite{Carroll2004}.

Berdasarkan elemen garis Schwarzschild pada Persamaan~\eqref{eq:metrik-schwarzschild}, komponen radial kovarian tensor metrik diberikan oleh $g_{rr}=f(r)^{-1}$. Karena metrik Schwarzschild bersifat diagonal dan $g^{\mu\nu}$ merupakan invers dari $g_{\mu\nu}$, komponen radial kontravariannya dapat dituliskan sebagai
\begin{equation}
g^{rr}=f(r)=1-\frac{2M}{r}.
\label{eq:grr-invers-schwarzschild}
\end{equation}
Kriteria $g^{rr}=0$ ini konsisten dengan penutupan kerucut cahaya yang telah ditinjau pada Persamaan~\eqref{eq:laju-cahaya-radial}, karena keduanya bergantung pada syarat $f(r) \to 0$ pada radius yang sama. Letak horizon peristiwa ditentukan dengan mencari akar fungsi metrik tersebut, sehingga diperoleh
\begin{equation}
f(r_h)=1-\frac{2M}{r_h}=0 \quad \Rightarrow \quad r_h=2M.
\label{eq:horizon-schwarzschild}
\end{equation}
Nilai $r_h=2M$ dikenal sebagai radius Schwarzschild dan menandai letak horizon peristiwa lubang hitam Schwarzschild \parencite{Carroll2004}.

Permukaan $r = 2M$ yang diperoleh dari kriteria $g^{rr} = 0$ ini bukanlah singularitas 
fisis, melainkan singularitas koordinat yang muncul akibat pemilihan sistem koordinat 
Schwarzschild dan dapat dihilangkan dengan pemilihan koordinat yang sesuai. Adapun 
singularitas yang sesungguhnya terletak pada $r = 0$, tempat kelengkungan ruang-waktu 
menjadi tak hingga dan tidak dapat dihilangkan oleh transformasi koordinat apa pun 
\parencite{Frolov1998}. Dengan demikian, $r = 2M$ merupakan \textit{null hypersurface} 
yang menandai batas kausal sebagaimana telah didefinisikan, sedangkan $r = 0$ merupakan 
singularitas fisis yang sesungguhnya berada di pusatnya.

Pada jarak yang jauh dari horizon peristiwa ($r \to \infty$), fungsi metrik mendekati 
nilai satu sehingga ruang-waktu Schwarzschild bersifat datar secara asimtotik. Sifat 
datar asimtotik ini menentukan bagaimana kondisi batas diterapkan pada daerah yang jauh 
dari lubang hitam, sebagaimana akan ditinjau pada pembahasan perturbasi medan skalar.

\subsection{Lubang Hitam Schwarzschild Anti-de Sitter}
\label{subsec:lh-sads}

\noindent Lubang hitam Schwarzschild Anti-de Sitter merupakan bentuk generalisasi lubang hitam
Schwarzschild pada ruang-waktu dengan konstanta Kosmologi bernilai negatif. Geometri
ruang-waktunya dideskripsikan oleh fungsi metrik
\begin{equation}
    f(r) = 1 - \frac{2M}{r} - \frac{\Lambda}{3} r^2, \qquad \Lambda < 0.
    \label{eq:f-sads}
\end{equation}
Dengan mendefinisikan jari-jari anti-de Sitter $L$ melalui hubungan $\Lambda = -3/L^2$,
fungsi metrik dapat dituliskan kembali sebagai
\begin{equation}
    f(r) = 1 - \frac{2M}{r} + \frac{r^2}{L^2}.
    \label{eq:f-sads-L}
\end{equation}

Sebagaimana pada kasus Schwarzschild, letak horizon peristiwa dapat ditentukan dari akar
fungsi metrik yang bernilai nol. Akan tetapi, kehadiran suku tambahan pada Persamaan
\eqref{eq:f-sads-L} menyebabkan persamaan tidak lagi memiliki penyelesaian sesederhana
pada latar belakang Schwarzschild, sehingga horizon peristiwa $r_+$ didefinisikan sebagai
akar positif terbesar dari persamaan
\begin{equation}
    f(r_+) = 1 - \frac{2M}{r_+} + \frac{r_+^2}{L^2} = 0.
    \label{eq:horizon-sads}
\end{equation}
Dengan mengalikan persamaan tersebut dengan $r_+ L^2$, syarat horizon dapat dinyatakan
dalam bentuk polinomial
\begin{equation}
    r_+^3 + L^2 r_+ - 2M L^2 = 0,
    \label{eq:horizon-sads-polinomial}
\end{equation}
yang merupakan persamaan pangkat tiga.

Karena Persamaan \eqref{eq:horizon-sads-polinomial} telah berbentuk tereduksi, akar-akarnya
dapat diperoleh melalui rumus Cardano. Akar real dari persamaan tersebut adalah
\begin{equation}
    r_+ = \sqrt[3]{M L^2 + \sqrt{M^2 L^4 + \frac{L^6}{27}}}
        + \sqrt[3]{M L^2 - \sqrt{M^2 L^4 + \frac{L^6}{27}}},
    \label{eq:horizon-cardano-sads}
\end{equation}
sedangkan kedua akar sisanya berbentuk pasangan kompleks konjugat
\begin{equation}
    r_{2,3} = -\frac{r_+}{2} \pm i\,\sqrt{L^2 + \frac{3}{4} r_+^2}.
    \label{eq:akar-kompleks-sads}
\end{equation}
Dengan demikian, Persamaan \eqref{eq:horizon-sads-polinomial} memiliki tepat satu akar real
positif yang diidentifikasi sebagai horizon peristiwa $r_+$, serta sepasang akar kompleks
konjugat berbagian real negatif yang tidak bersesuaian dengan jari-jari fisis mana pun
\parencite{Horowitz2000}.

Perbedaan mendasar dari lubang hitam Schwarzschild terletak pada perilaku geometri di daerah
yang sangat jauh dari horizon peristiwa. Pada ruang-waktu Schwarzschild Anti-de Sitter (SAdS),
suku tambahan pada fungsi metrik tumbuh tanpa batas ketika $r \rightarrow \infty$, sehingga
ruang-waktu tidak lagi mendekati ruang datar, melainkan mendekati geometri Anti-de Sitter
yang memiliki kelengkungan negatif konstan \parencite{Horowitz2000}. Akibatnya, struktur
ruang-waktu pada daerah asimtotik menjadi sangat berbeda dengan ruang-waktu Schwarzschild.

Meskipun daerah tersebut berada pada jarak tak hingga, geometri di sana masih dapat dipelajari
melalui suatu transformasi konformal yang memetakan titik-titik di tak hingga menjadi suatu
batas yang berhingga. Batas hasil pemetaan ini dikenal sebagai batas konformal. Dengan
demikian, sifat geometri pada daerah yang sangat jauh dari lubang hitam tetap dapat dianalisis
tanpa harus berhadapan secara langsung dengan titik yang berada di tak hingga.

Karakteristik batas konformal bergantung pada jenis ruang-waktu yang ditinjau. Pada
ruang-waktu Schwarzschild yang bersifat datar secara asimtotik ($\Lambda = 0$), batas
konformal bersifat \textit{light-like}. Hal ini menunjukkan bahwa cahaya maupun gelombang
yang merambat menjauhi lubang hitam dapat terus bergerak menuju tak hingga tanpa kembali
lagi ke daerah sekitar lubang hitam. Sebaliknya, pada ruang-waktu Anti-de Sitter
($\Lambda < 0$), batas konformal bersifat \textit{timelike}. Secara fisis, batas ini dapat
dibayangkan menyerupai sebuah dinding pemantul yang mengurung ruang-waktu, sehingga gelombang
yang merambat ke arah luar tidak dapat terus menjalar menuju tak hingga, melainkan akan
dipantulkan kembali.